\chapter{Colinéarité, déterminant et bases}
\section*{Exercice 9.1 -- Déterminants}
Calculer les déterminants de $(2,5)$ et $(6,15)$, puis de $(2,5)$ et $(3,7)$. Conclure.

\section*{Exercice 9.2 -- Alignement}
$A(1,-1)$, $B(3,2)$, $C(5,5)$. Démontrer l'alignement.

\section*{Exercice 9.3 -- Paramètre}
Déterminer $t$ pour que $(t,6)$ et $(2,3)$ soient colinéaires.

\section*{Exercice 9.4 -- Coordonnée nulle}
Expliquer pourquoi $(0,4)$ et $(0,-7)$ sont colinéaires sans écrire de quotient de dénominateur nul.

\section*{Exercice 9.5 -- Dans une autre base}
Exprimer $\vec w=(7,1)$ dans la base $\vec u=(1,2)$, $\vec v=(2,-1)$.

\section*{Exercice 9.6 -- Parallèles ou confondues}
$A(0,0)$, $B(2,1)$, $C(0,3)$, $D(4,5)$. Comparer les droites $(AB)$ et $(CD)$.
